\documentclass[tikz,border=7pt]{standalone}
% flashcards-title: Not-to-scale segment diagram
% flashcards-desc: A visually longer segment is labeled 4 centimeters with one tick, while a visually shorter segment is labeled 6 centimeters with two ticks.
\input{figures/tikz-style.tex}
\begin{document}
\begin{tikzpicture}[x=1cm,y=1cm]
  \node[geom panel,minimum width=8.8cm,minimum height=4.9cm] at (0,0) {};
  \node[geom label] at (0,1.85) {diagram not drawn to scale};
  \draw[geom line] (-3.4,0.85)--(2.9,0.85);
  \node[geom point] at (-3.4,0.85) {};
  \node[geom point] at (2.9,0.85) {};
  \node[geom small,fill=white,inner sep=2pt] at (-0.25,0.85) {$4\ \mathrm{cm}$};
  \draw[geom line] (-2.0,-0.55)--(1.6,-0.55);
  \node[geom point] at (-2.0,-0.55) {};
  \node[geom point] at (1.6,-0.55) {};
  \node[geom small,fill=white,inner sep=2pt] at (-0.2,-0.55) {$6\ \mathrm{cm}$};
  \draw[geom line,line width=1.2pt] (-0.42,1.04)--(-0.12,0.66);
  \draw[geom line,line width=1.2pt] (-0.48,-0.36)--(-0.18,-0.74);
  \draw[geom line,line width=1.2pt] (-0.22,-0.36)--(0.08,-0.74);
  \node[geom small,align=center] at (0,-1.65) {values and matching marks state the mathematical information};
\end{tikzpicture}
\end{document}
